\chapter*{理论脉络}
\addcontentsline{toc}{chapter}{理论脉络}

各章按照理论依赖关系单向展开。连续谱、投影与对称性先建立相对论经典场的数学语言；粒子产生与湮灭随后要求进入 Fock 空间和场相关函数。QED、QCD 与标准模型给出这些结构在基本相互作用中的实现，自由电子与光、介质及测量问题再建立在这些母理论之上。具体装置、材料与低维模型最后作为一般方程的参数化或受控极限出现。

\begin{longtable}{p{0.15\textwidth}p{0.34\textwidth}p{0.43\textwidth}}
\toprule
章节 & 核心问题 & 后续作用\\
\midrule
\chapref{chap:math-foundations} & 连续变量怎样作 Fourier 分析；怎样处理分布、连续谱、投影消元、二次积分和连续对称变换。 & 固定全书共同的数学运算、归一化和误差控制语言。\\
\chapref{chap:classical-fields} & 哪些线性变换保持相对论时空间隔；这些变换怎样作用在不同类型的相对论场变量上；电磁场怎样由局域作用量和边界条件决定。 & 给电子相对论动力学、经典电磁场以及后续量子化提供同一套时空与符号约定。\\
\chapref{chap:qft-foundations} & 当粒子可以产生和湮灭时，怎样建立一个允许粒子数改变的 Hilbert 空间；局域场怎样量子化；相关函数怎样连接到散射实验。 & 为所有后续粒子与光子过程提供共同的量子多粒子框架。\\
\chapref{chap:qed} & 把电子场和电磁场同时量子化以后，怎样系统计算电子--光子散射、辐射和量子修正。 & 给自由电子与量子光场相互作用提供微观母理论。\\
\chapref{chap:qcd} & 若内部“荷”有多个分量且不同变换不交换，规范场自身为什么也会相互作用；强相互作用怎样随能标变化。 & 展示非交换内部对称的完整动力学，并为粒子物理统一结构准备颜色相互作用部分。\\
\chapref{chap:standard-model} & 怎样把强、弱和电磁相互作用的场内容放进同一局域理论，并解释规范玻色子与费米子的质量、味混合和低能极限。 & 固定后续若涉及基本粒子过程时所调用的场内容和耦合。\\
\chapref{chap:general-interaction} & 为什么自由电子首先是连续正能质量壳，而不是一串离散能级；怎样从 QED 投影得到“电子 + 任意环境”的时域母 Hamiltonian，并判断被消去通道何时留下记忆或损耗。 & 给出自由电子量子光学最上层的连续谱母理论，使正常模、Green 张量、单模和低维模型都有明确来源。\\
\chapref{chap:external-pinem}至\chapref{chap:periodic-recoil} & 精确质量壳失谐怎样选择有限时间交换通道；在什么条件下真实反冲格可以进一步约化成 PINEM/QPINEM 均匀能量梯或少态共振。 & 把“边带”“反冲”“量子交换”统一为同一连续电子谱的不同受控尺度极限。\\
\chapref{chap:prop-measure}至\chapref{chap:media-radiation} & 相互作用后产生的相干怎样经精确色散传播、环境约化和测量变成数据；量子跃迁流何时才可进一步替换为经典给定电流。 & 把联合量子态和材料 Green 张量最终连接到能谱、成像、计数与参数反演。\\
\bottomrule
\end{longtable}

各章沿表中的依赖关系顺次展开。新对象第一次进入正文时，先由它所解决的物理问题引出，再给出数学定义，因此阅读时无需预先掌握后续章节的专有名词。
